\section{Treap impl\'icito}
\whenuse{Secuencia cuyo largo u orden cambia: insertar/borrar en el medio,
invertir, cortar o consultar segmentos. Para un arreglo de largo fijo suele ser
m\'as simple un Segment Tree.}

El recorrido inorder del \'arbol es la secuencia. No guarda claves: el tama\~no
de cada sub\'arbol determina la posici\'on. \texttt{split(t,k)} separa los
primeros $k$ elementos del resto y \texttt{merge(a,b)} concatena dos secuencias.
Para operar sobre $[l,r)$ se lo a\'isla con dos splits y luego se vuelven a unir
las tres partes.

\complexity{split, merge, inserci\'on, borrado, inversi\'on y consulta
$O(\log n)$ esperado; memoria $O(n)$.}

\begin{lstlisting}[style=cpp]
struct ImplicitTreap {
    struct Node {
        ll value, sum;
        unsigned priority;
        int size = 1;
        bool reversed = false;
        Node *left_child = nullptr, *right_child = nullptr;
        Node(ll initial_value, unsigned initial_priority)
            : value(initial_value), sum(initial_value),
              priority(initial_priority) {}
    };

    Node *root = nullptr;

    static unsigned next_priority() {
        static mt19937 generator(chrono::steady_clock::now()
                                 .time_since_epoch().count());
        return generator();
    }

    static int node_size(Node *node) { return node ? node->size : 0; }
    static ll node_sum(Node *node) { return node ? node->sum : 0; }
    static void push(Node *node) {
        if (!node || !node->reversed) return;
        swap(node->left_child, node->right_child);
        if (node->left_child) node->left_child->reversed ^= true;
        if (node->right_child) node->right_child->reversed ^= true;
        node->reversed = false;
    }
    static void pull(Node *node) {
        node->size = 1 + node_size(node->left_child)
                       + node_size(node->right_child);
        node->sum = node->value + node_sum(node->left_child)
                                + node_sum(node->right_child);
    }

    static void split(Node *node, int left_size,
                      Node *&left_tree, Node *&right_tree) {
        if (!node) { left_tree = right_tree = nullptr; return; }
        push(node);
        if (node_size(node->left_child) >= left_size) {
            split(node->left_child, left_size,
                  left_tree, node->left_child);
            right_tree = node; pull(right_tree);
        } else {
            int remaining = left_size - node_size(node->left_child) - 1;
            split(node->right_child, remaining,
                  node->right_child, right_tree);
            left_tree = node; pull(left_tree);
        }
    }

    static Node *merge(Node *left_tree, Node *right_tree) {
        if (!left_tree || !right_tree)
            return left_tree ? left_tree : right_tree;
        push(left_tree); push(right_tree);
        if (left_tree->priority > right_tree->priority) {
            left_tree->right_child = merge(left_tree->right_child, right_tree);
            pull(left_tree); return left_tree;
        }
        right_tree->left_child = merge(left_tree, right_tree->left_child);
        pull(right_tree); return right_tree;
    }

    void insert(int position, ll value) {
        assert(0 <= position && position <= node_size(root));
        Node *left, *right;
        split(root, position, left, right);
        Node *node = new Node(value, next_priority());
        root = merge(merge(left, node), right);
    }

    void erase(int position) {
        assert(0 <= position && position < node_size(root));
        Node *left, *middle, *right;
        split(root, position, left, middle);
        split(middle, 1, middle, right);
        delete middle;
        root = merge(left, right);
    }

    void reverse_range(int left, int right) { // [left,right)
        assert(0 <= left && left <= right && right <= node_size(root));
        Node *prefix, *middle, *suffix;
        split(root, right, middle, suffix);
        split(middle, left, prefix, middle);
        if (middle) middle->reversed ^= true;
        root = merge(merge(prefix, middle), suffix);
    }

    ll range_sum(int left, int right) { // [left,right)
        assert(0 <= left && left <= right && right <= node_size(root));
        Node *prefix, *middle, *suffix;
        split(root, right, middle, suffix);
        split(middle, left, prefix, middle);
        ll answer = node_sum(middle);
        root = merge(merge(prefix, middle), suffix);
        return answer;
    }
};
\end{lstlisting}

\subsection{Ejemplo de uso}

\begin{lstlisting}[style=cpp]
void implicit_treap_example() {
    ImplicitTreap sequence;
    sequence.insert(0, 10); // [10]
    sequence.insert(1, 20); // [10,20]
    sequence.insert(2, 30); // [10,20,30]
    sequence.insert(1, 5);  // [10,5,20,30]

    cout << sequence.range_sum(1, 4) << '\n'; // 55
    sequence.reverse_range(0, 3);              // [20,5,10,30]
    cout << sequence.range_sum(0, 2) << '\n'; // 25
    sequence.erase(1);                         // [20,10,30]
}
\end{lstlisting}

La prioridad aleatoria mantiene el \'arbol balanceado en esperanza. La bandera
\texttt{reversed} es lazy: invertir un segmento s\'olo marca su ra\'iz y el
cambio se propaga cuando se baja por ese sub\'arbol.

\warning{Esta versi\'on reserva nodos con \texttt{new}. Si se crean y destruyen
muchas secuencias, agregar un destructor o usar un pool de memoria. No copiar la
estructura directamente: dos copias compartir\'ian los mismos punteros.}
